\section{Introduction}\label{sec:introduction}

The triangular Hilbert form couples three functions of two variables,
each pair sharing one variable. We study its dyadic model, in which the
intervals are constrained by binary addition without carries. The
estimate below imposes no structural condition on any of the three
inputs.

Write \(\R_+=[0,\infty)\). For \(k\in\Z\), let
\[
 \cD_k=\{[2^kn,2^k(n+1)):n\in\Z_{\ge0}\}.
\]
For \(I\in\cD_k\), denote its index by \(n_I\) and set
\[
 h_I=\mathbf1_{I_{\mathrm{left}}}-\mathbf1_{I_{\mathrm{right}}}.
\]
Let \(\oplus\) denote bitwise XOR of nonnegative integers, and define
the set of admissible triples
\[
 \cA_k=\{(I_0,I_1,I_2)\in\cD_k^3:
                  n_{I_0}\oplus n_{I_1}\oplus n_{I_2}=0\}.
\]
For \(\mathbf I\in\cA_k\) and three real functions \(F_v\) on
\(\R_+^2\), put
\begin{equation}\label{eq:local-form}
 L_{\mathbf I}(F_0,F_1,F_2)
 =2^{-k}\int_{I_0\times I_1\times I_2}
 F_0(x_0,x_1)F_1(x_1,x_2)F_2(x_2,x_0)
 \prod_{v=0}^2 h_{I_v}(x_v)\dd x_0\dd x_1\dd x_2.
\end{equation}
The dyadic triangular Hilbert form with a finite set of scales
\(S\subset\Z\) and real coefficients \(\abs{\epsilon_k}\le1\) is
\begin{equation}\label{eq:given-form}
 \Lambda^{\mathrm d}_{S,\epsilon}(F_0,F_1,F_2)
 =\sum_{k\in S}\epsilon_k\sum_{\mathbf I\in\cA_k}
                  L_{\mathbf I}(F_0,F_1,F_2).
\end{equation}

\begin{theorem}\label{thm:main}
For every finite \(S\subset\Z\) and all real bounded compactly
supported measurable functions \(F_0,F_1,F_2\) on \(\R_+^2\),
\begin{equation}\label{eq:main}
 \sum_{k\in S}\sum_{\mathbf I\in\cA_k}
       \abs{L_{\mathbf I}(F_0,F_1,F_2)}
 \le 40\prod_{v=0}^2\norm{F_v}_{L^3(\R_+^2)}.
\end{equation}
\end{theorem}

In particular, \eqref{eq:given-form} satisfies the same product bound,
uniformly in the scales and coefficients. The absolute values in
\eqref{eq:main} also allow a separate coefficient of modulus at most one
for each admissible triple. The estimate extends to all real \(L^3\)
inputs by density, as explained in Section~\ref{sec:summation}.

\subsection{Context and relation to earlier work}
The continuous triangular Hilbert form has the same cyclic arrangement
of inputs as \eqref{eq:local-form}, integrated over \(\R^3\) against the
principal-value kernel \(1/(x_0+x_1+x_2)\). Its symmetric \(L^3\) estimate is formulated by
Thiele in \cite[Section~9, Problem~13]{ThieleProblem2017}. That
contribution also displays the precise dyadic model
\eqref{eq:given-form}, with one coefficient for each scale, and asks
for the analogous bound. Theorem~\ref{thm:main} answers this dyadic
question affirmatively and gives the stronger sum of absolute local
contributions. The continuous form and its dyadic model are distinct
operators; the theorem here is a statement about the latter.

Kova\v{c}, Thiele and Zorin-Kranich introduced a dyadic triangular
model with coefficients allowed to vary among interval triples
\cite[Equations~(1.5)--(1.6)]{KTZ2015}. Their Theorem~1.7 gives
uniform estimates when one input has either a specified dependence
on a Walsh linear combination of its coordinates or a specified
fiberwise character structure. Those cases already encompass dyadic
versions of Carleson's operator and uniform bounds for the bilinear
Hilbert transform, as explained in their Appendix~A. Our estimate
allows all three inputs to vary freely. The cyclic operator-trace
representation used below also appears in their Equation~(1.5).

Another line of work measures cancellation through the growth of
truncated norms. Tao obtained sublogarithmic growth for
one-dimensional multilinear Hilbert transforms
\cite[Theorem~1.2]{Tao2016}; Zorin-Kranich proved the corresponding
sublinear growth in the number of smooth scales for simplex Hilbert
transforms \cite[Theorem~1.3]{ZK2017}. Durcik, Kova\v{c} and Thiele
then obtained power-type cancellation in both the continuous and
dyadic settings \cite[Theorems~1 and~4]{DKT2019}. For the triangular
dyadic form, their bound over \(m\) consecutive scales has a factor
\(m^{1/2}\) at exponents \((4,4,2)\); cyclic permutation and multilinear
interpolation give the same growth at \((3,3,3)\). Theorem~\ref{thm:main}
removes this scale dependence at the symmetric tuple.

The strategy of controlling an entangled form by the change of a
nonnegative energy has precedents in Kova\v{c}'s Bellman-function
framework \cite[Section~2]{Kovac2011} and his telescoping identity
for the twisted paraproduct \cite[Lemma~2]{Kovac2012}.
These provide a methodological precedent, rather than a bound for the
present triangular form. In particular, Kova\v{c} explicitly identifies
the obstruction to applying that framework to non-bipartite graphs,
including the triangle, in \cite[Section~6.2]{Kovac2011}.
Here the energy is a trace of a matrix power, and the averaging is
over the four children selected by the XOR constraint.

The matrix energy and the dimension-independent trace inequality are
taken from the companion manuscript \cite[Sections~2--3, Lemmas~2.1,
2.2 and~3.1]{OpenAI2026}; we reproduce their real finite-dimensional
proofs. Their underlying matrix calculus is classical: the
square-root derivative is a special case of the divided-difference
formula in \cite[Section~1, Equation~(1)]{Daletskii1957}, and the order
comparison uses monotonicity of the positive square root
\cite[Theorem, p.~309]{Pedersen1972}.


\subsection{The energy argument}

The proof reduces the form to finite matrices by taking inputs constant
on sufficiently small dyadic squares. On a fixed admissible triple,
each input becomes a matrix between the atom spaces of its two
variables. At each variable, we join its two incident matrices into a
single row operator \(R\) and assign it the energy
\(\tr((R^*R)^{3/2})\). Let \(e(\mathbf I)\) be the sum of these three
energies, and let \(\operatorname{ch}(\mathbf I)\) consist of the four
admissible triples obtained by halving each interval. The key estimate
is
\begin{equation}\label{eq:overview-gap}
 \abs{L_{\mathbf I}(F_0,F_1,F_2)}
 \le C_0\,2^{2k}
 \left(\frac14\sum_{\mathbf I'\in\operatorname{ch}(\mathbf I)}
                 e(\mathbf I')-e(\mathbf I)\right),
 \qquad \mathbf I\in\cA_k,
\end{equation}
with an absolute constant \(C_0\). At each scale, the weighted total
energy \(\sum_{\mathbf I\in\cA_k}2^{2k}e(\mathbf I)\) is bounded by an
absolute constant times the sum of the cubed input norms.

To prove \eqref{eq:overview-gap}, a matrix trace estimate bounds
the triangular product by three Hessian quadratic forms of
\(P\mapsto\tr(P^{3/2})\). Conditional on one child half, each neighboring
half remains equally likely to be left or right. The two neighbors are
correlated, but \(RR^*\)
is a sum of the two edge Gram matrices; it contains no product between
their multipliers and has the same nonzero spectrum as \(R^*R\).
Convexity therefore converts the Hessian bound into an increase of the
child energy. Writing the local integrals in normalized interval
measures leaves the weight \(2^{2k}\). Each child has one parent and
weight \(2^{2(k-1)}\), one quarter of the parent's weight, so these
increases telescope.

Section~\ref{sec:matrix} establishes the energy
inequalities, Section~\ref{sec:trace} proves the trace estimate,
Section~\ref{sec:dyadic} supplies the conditional-averaging argument,
and Section~\ref{sec:summation} sums the resulting differences and
passes to measurable inputs.
